\documentclass[12pt]{article}

\usepackage[margin=1in]{geometry}
\usepackage{amsmath,amssymb,amsthm}
\usepackage{booktabs}
\usepackage{titlesec}
\titleformat{\section}{\normalfont\bfseries}{\thesection.}{1em}{\MakeUppercase}
\titleformat{\subsection}{\normalfont\bfseries}{\thesubsection}{1em}{}
\usepackage{caption}
\captionsetup{labelsep=period,justification=raggedright,singlelinecheck=false}
\usepackage{natbib}
\usepackage[hidelinks]{hyperref}
\usepackage{setspace}
\usepackage{pgfplots}
\pgfplotsset{compat=1.18}

\bibpunct{(}{)}{;}{a}{}{,}

% Allow extra stretch in math-heavy paragraphs rather than overflowing the margin.
\emergencystretch=2em

% Let floats occupy more of a page before being pushed to a page of their own.
\renewcommand{\topfraction}{0.9}
\renewcommand{\bottomfraction}{0.8}
\renewcommand{\textfraction}{0.07}
\renewcommand{\floatpagefraction}{0.75}

\newtheorem{theorem}{Theorem}
\newtheorem{proposition}{Proposition}
\newtheorem{lemma}{Lemma}
\newtheorem{corollary}{Corollary}
\theoremstyle{definition}
\newtheorem{assumption}{Assumption}
\newtheorem{definition}{Definition}

\numberwithin{equation}{section}

% Captions sit above tables, and the article class leaves no space below a
% caption, so a caption running to two lines collides with \toprule.
\setlength{\abovecaptionskip}{10pt}
\setlength{\belowcaptionskip}{7pt}

\title{What Pre-Trends Cannot See:\\[2pt]
Testing for Treatment-Induced Recording Change\\[2pt]
with Few Clusters}

\author{}

\date{November 9, 2026}

\begin{document}
\maketitle

\begin{abstract}
\noindent
A policy that changes how an outcome is recorded is invisible to pre-treatment
diagnostics. This is the position of any difference-in-differences design whose
outcome is drawn from an administrative record that the policy, or a reform
arriving with it, also rebuilt. The parallel trends deviation reproducing a recording change is
exactly zero before treatment, so pre-trend tests reject no more often than
their nominal level and sensitivity analysis calibrated from the pre-period loses coverage. Falsification tests across outcome categories graded by
recording discretion supply the missing information. Run as one joint test with
cluster-robust standard errors, the chi-square calibration is oversized: with
51 clusters, 13 restrictions and 5 categories its null rejection rate is 0.389
against a nominal 0.05. We give three randomization calibrations and
characterize the size of each: label permutation, exact in finite samples under
category exchangeability and valid for any number of clusters; cluster sign
changes, exact under sign symmetry of the cluster contributions; and the
restricted wild cluster bootstrap. Label randomization carries a resolution
floor that requires at least five categories at the five percent level. Whether
a design can detect the recording change its institution implies is computable
before the outcome is examined, and it turns on alignment as well as power: the
magnitude that matters is the projection of the institutional response on the
exposure index, and in all three panels we screen it falls an order of magnitude
short of what the design detects.
\end{abstract}

\smallskip
\noindent\textbf{Keywords:} difference-in-differences; measurement error; randomization inference; sensitivity analysis; wild bootstrap.


\smallskip
\noindent\textbf{JEL codes:} C12, C15, C18, C23, C81.

\newpage
\doublespacing

\section{Introduction}
\label{sec:intro}

Applied researchers routinely assess an identifying assumption by asking whether
an estimated treatment effect behaves across outcomes the way the assumption
predicts, choosing outcomes that differ in their exposure to a suspected
violation. For example, \citet{iyer2012political} separate reporting from
incidence by contrasting documented crimes against women with crimes against men
and gender neutral crimes, and \citet{miller2019female} run the same contrast for
violent victimization by sex of the victim. \citet{weisburd2024crs} grade crime
categories by their reporting propensity and compare estimated effects across the
grades. \citet{baicker2013oregon} report effects of Medicaid coverage on measured
blood pressure and cholesterol alongside self reported diagnoses of the
corresponding conditions. In epidemiology the same device is formalized as the
negative control outcome \citep{lipsitch2010negative}, and
\citet{nunez2025detection} grade candidate controls by how much of the
ascertainment process they share with the outcome of interest, which is the
cross section this paper uses; \citet{zhang2025nco} calibrate a
difference-in-differences estimate with negative controls, though against
unmeasured confounding rather than a recording response. Run as a single test, the
comparison becomes a joint restriction on interaction coefficients, and its $p$
value is read from a chi-square reference distribution with cluster-robust standard errors.

In this paper we study the case in which the suspected violation is a change in
the way the outcome is recorded rather than a change in the outcome itself.
\citet{campbell1963experimental} named this the instrumentation threat and
advised against shifting the measuring instrument and the policy at the same
time, advice that is of little use once the policy is what moved the instrument.

How much a recording layer can move is not a small quantity, and it has been
measured where the recording technology was assigned at random.
\citet{singh2024administrative} tests over $2{,}400$ schools in Andhra Pradesh
and finds that paper-based assessments return scores $16$ to $20$ percentage
points above an independently proctored retest of the same students, while
tablet-based assessments match the retest. The skill being measured is the same
in both arms and the instrument is what differs, which is a recording response
in the sense of Assumption \ref{as:rec} with the wedge identified by design
rather than assumed. That such a change should coincide with a policy is common
rather than exotic wherever administrative recording is being rebuilt, which is
much of the setting in which quantitative work on developing and emerging
economies is done \citep{singh2020myths}. Where the reform to the record arrives
with the policy, Assumption \ref{as:rec} is the data generating process and not
a hypothetical, and the analyst who reaches for an event-study plot is reaching
for the one diagnostic that Corollary \ref{cor:blind} says cannot help.
We show that for any recording response there exists a structure
with no recording response whatever, and with a violation of parallel trends that
is exactly zero throughout the pre-treatment period, which generates the same
distribution of the observed panel cell by cell. Because the two structures are
indistinguishable in the data, no diagnostic computed from the observed panel and
the cohort labels rejects more often than its nominal level, whatever
the magnitude of the recording response, and pre-trend tests, placebo in time
tests and leads and lags plots are all functions of exactly that information.

The same construction tells us what a sensitivity analysis can deliver. The
standard response to an untestable assumption is not to test it but to report a
confidence set that is honest over a class of deviations
\citep{rambachan2023credible}. In the relative magnitudes class the bound on post-treatment deviations is a fixed multiple of the largest pre-treatment deviation.
Under a recording response the quantity being scaled is identically zero, so the
class collapses to a single point for every value of its sensitivity parameter
and the procedure reports no breakdown value at all. In simulations calibrated to
a latent effect of one and a recording wedge of $0.6$, and computed with a
reference implementation rather than by hand, coverage of the latent effect at a
sensitivity parameter of one falls from $0.955$ at $N = 400$ to $0.008$ at
$N = 6{,}400$ and to zero thereafter. An interval whose bound is taken from the
recording institution rather than from the pre-period covers at one throughout.
More data does not repair the pattern but is instead what produces it, because
the calibration input falls at the sampling rate while the wedge it is meant to
bound does not move at all.

Information about a recording response has to come from somewhere other than the
pre period. We take it from the cross section of outcome categories, which differ
in how much discretion the recording system has over them. That cross sectional
comparison is the falsification test with which we began, and the remainder of
the paper is concerned with how to calibrate it when the number of clusters is
small, which is the situation in which these tests are almost always run. The
calibration in current use fails badly at exactly those sample sizes. At $51$
clusters, $13$ restrictions and $5$ categories we measure a null rejection rate
of $0.389$ against a nominal $0.05$, and at $20$ clusters the rate is $0.905$.

Why it fails has a closed form, and the closed form is what turns the
failure into a design rule. When every cluster carries the same design
weight, the cluster robust Wald statistic is exactly $\omega_2$ times a
Hotelling $T^2$ with dimension equal to the number of restrictions and
degrees of freedom equal to one less than the number of clusters
\citep{hotelling1931}. The chi-square reference ignores the estimation of
the meat matrix from the cluster scores, and the identity prices that
omission. Equal weights are what make the statement exact and they hold in
no design we run, including both applications, so Proposition
\ref{prop:anatomy} earns its place as a diagnostic and not as a
calibration we recommend; its own reference distribution rejects between
$0.082$ and $0.175$ across the twelve cells of Table \ref{tab:size}. What
survives the concession is the rule. The distortion is governed jointly by
the number of restrictions and the number of clusters, so the chi-square
calibration is valid only when the first grows more slowly than the two
thirds power of the second, and a researcher can check that before
collecting an outcome. At $36.4$ effective clusters size stays below
$0.075$ only when two or fewer restrictions are imposed.

We then give three calibrations that do not require equal design weights, and we
characterize the size of each. The first permutes the exposure labels across
outcome categories, and we show that under an exchangeability condition on the
joint law of the category data arrays this test has size at most $\alpha$ in
finite samples for every number of units, periods, clusters and categories, with
no distributional assumption and no condition whatever on the number of
clusters. Its randomization group is the set of relabelings, whose order bounds
the resolution of the test, and the smallest attainable $p$ value is $2/C!$ when the statistic is
even in the labels and the exposure indices are symmetric about their mean, and
$1/C!$ otherwise, so a two-sided test at the five percent level with equally
spaced indices requires at least five categories. The second procedure randomizes
sign changes over clusters, and it is exactly level $\alpha$ when the cluster
contributions are sign symmetric under the null and asymptotically level $\alpha$
when they are not. The third is the restricted wild cluster bootstrap of
\citet{cameron2008bootstrap}, for which we report simulated size rather than a
theorem. Both cluster level procedures remain usable when there are too few
categories for permutation.

Two further results concern power rather than size. We show that the Wald
statistic is not optimal against the alternative these tests are designed to
detect, and that when the shape of the response across horizons is known up to
scale a directed statistic is uniformly most powerful invariant, which puts the
ratio of minimum detectable responses at power $0.80$ at $0.393$ at $36.4$
effective clusters. Label randomization carries a price of its own, requiring
$1.35$ times the signal at five categories and $1.15$ at eight relative to a
normal benchmark. We use the label permutation test as our finite-sample benchmark, because its reference distribution is valid for every number of
clusters and because it is the only one of the three that holds its size across
every design we try. The directed cluster sign-change test is the more powerful
procedure and is the one to use when each cluster is a single unit, which is the
case in both applications here, but its null rejection rate rises to $0.094$
when clusters hold five correlated units and to $0.086$ when their sizes are
unequal. Its two-sided counterpart is far less affected, at $0.057$ and
$0.063$, so a cluster holding several units is a reason to give up the power of
the directed statistic rather than the cluster level calibration.

A falsification test is worth running only against responses it could
detect, and whether it can is computable before any outcome is examined. We
give that calculation and run it over forty eight specifications of each panel.
No specification of either panel is adequate, and the reason in the crime panel
is not a shortage of data. The response the move to incident-based reporting
implies is real, averaging $0.011$ log points across index categories, but it
acts through police counting rules while the exposure index grades victim
reporting propensity, so the two are nearly orthogonal and the projection of the
one on the other is $0.0044$ against a minimum detectable response of $0.052$.
A design can therefore fail this screen through misalignment rather than through
power, and which of the two it is can be read off before the outcome is
examined.

The rest of the paper is organized as follows. Section \ref{sec:why} studies what a pre-trend test and a pre-trend calibrated sensitivity analysis can reveal about a recording response. Section \ref{sec:setup} states the design and the null. Sections \ref{sec:anatomy} to \ref{sec:power} give the exact null law, the three calibrations and the power results. Section \ref{sec:sim} reports size and power, Section \ref{sec:designcalc} screens a design before the outcome is examined, and Section \ref{sec:app} applies the procedures to two panels. Proofs and further evidence are in the online Appendix.

\section{Related work}

Our exact null law concerns the cluster robust Wald statistic, which places it
in the few cluster literature of \citet{cameron2008bootstrap},
\citet{mackinnon2017wild} and \citet{carter2017cluster}, and the finite-sample argument is the classical group invariance construction of
\citet{lehmann2005testing}. Proposition \ref{prop:anatomy} does not introduce
the Hotelling family that \citet{bester2011cluster}, \citet{bell2002bias},
\citet{tipton2015small} and \citet{pustejovsky2018small} develop. It states an
algebraic identity and, under Assumption \ref{as:I}, an exact law holding at
every cluster count and every number of restrictions with $G - m \ge 1$, rather
than at matched degrees of freedom or in a large cluster limit. The design that
motivates the problem, grading outcome categories by recording discretion, is
run by \citet{weisburd2024crs}; \citet{roth2022pretest} and
\citet{rambachan2023credible} treat falsification tests run as pre-trend tests.
The epidemiological statement of that design is
\citet{nunez2025detection}, who reach the same grading informally and record
what is missing from it: that low precision complicates reading a negative
control, and that no formal guidance exists on how large a response one has to
be able to detect. Sections \ref{sec:perm} to \ref{sec:designcalc} are an
answer to those two sentences.
The online Appendix gives the extended discussion.

\section{What the falsification test is for}
\label{sec:why}

That a policy can change how an outcome is recorded rather than the outcome
itself is not a new observation. It is \citet{campbell1963experimental}'s
\emph{instrumentation} threat to internal validity, and their illustration is
almost exactly the design studied below: ``the experimental change of putting
into office a new principal may produce a change in the recording of discipline
infractions rather than in the infraction rate itself.''
\citet{campbell1969reforms} makes the same point for policing, where a change in
administration reforms the bookkeeping alongside the policy, and we claim no part
of that observation for ourselves.

What we do claim is a correction to what the classical tradition concluded about
it. For the single series design, \citet{campbell1963experimental} mark
instrumentation as the one threat the pre-period cannot dispel, while maturation,
testing and regression are ruled out precisely because a pre series is available.
Then, for the multiple time series design, which adds a control series and is the
ancestor of the design used here, they write that ``the difference in the
selection of the groups operating in conjunction with maturation, instrumentation,
or regression, can hardly account for an apparent effect.'' Their reasoning concerns
instrument differences that arise from how the groups were selected, and a
recording response caused by the treatment is not a difference of that kind. No
control series can absorb it, because the control units never receive the
treatment that induces it. The distinction between an instrument difference
across groups and a recording response induced by the treatment is not drawn
there, and the applications we cite below suggest that it is not always drawn
now.

The consequence we draw is sharper than the classical warning. When the
contamination is a recording response, a pre-trend test and a sensitivity
analysis calibrated from pre trends are not merely low powered against it, since
neither carries any information about it at all. We state that claim precisely
in what follows, price it, and mark the boundary of what it covers.

The gap is a live one, and the online Appendix documents it: work that warns about differential recording and then offers an event-study as the remedy, work that reports a pre-trend test as evidence of parallel trends while conceding that the outcome may reflect disclosure rather than incidence, and work that attributes an undetected violation to low power. Corollary \ref{cor:blind} says that where the contamination is a recording response none of these is a matter of power. Where the threat is taken seriously the response is already the method proposed here: \citet{harasztosi2019minwage} settle a reporting story not from the pre-period but across outcomes, which is a cross outcome falsification test conducted informally, and \citet{geruso2020upcoding} obtain a placebo with power only because an institutional lag supplies a window in which the coding response cannot yet have operated. The rest of this paper supplies the reference distribution such a test needs when the number of clusters is small.

That observation generalizes into the one constructive diagnostic we can offer.
The equivalence in Lemma \ref{lem:oe} is exact because Assumption \ref{as:rec}
places the recording change at the treatment date. Where the recording system responds to
an announcement, a rule change or an administrative deadline that precedes
implementation, the wedge is nonzero before treatment and a placebo in time test
placed in that window has power against it. The same is true of
\citet{geruso2020upcoding}'s lag. We therefore recommend that a researcher who suspects a recording response look
for a period in which the policy was already known but the recording system had
not yet moved, and run the placebo there. Where no such period exists, no test
computed from the outcome panel can help, and that is the case the rest of this
paper addresses.

\subsection{A recording response}
\label{sec:recording}

Index units by $i$, periods by $t$, outcome categories by $c$, and adoption
cohorts by $g$, with $\ell_{it} = t - G_i$ the event time and $D_{it} =
\mathbf 1\{t \ge G_i\}$. Write $Y^{*}_{itc}$ for the latent outcome and
$Y_{itc}$ for what the recording system reports.

\begin{assumption}[Separable recording response]
\label{as:rec}
There exist $\mu_{ic}$, depending arbitrarily on $G_i$ and on the latent
outcomes, deterministic $\phi_{tc}$, and deterministic paths $m_c(\ell)$ with
$m_c(\ell) = 0$ for $\ell < 0$, such that
\begin{equation}
\label{eq:recmap}
Y_{itc} = Y^{*}_{itc} + \mu_{ic} + \phi_{tc} + m_c(\ell_{it})D_{it} + u_{itc},
\qquad E[u_{itc} \mid G_i] = 0 .
\end{equation}
\end{assumption}

Two features of Assumption \ref{as:rec} carry the results below and both are
restrictions. The distortion enters additively and separably, so a recording
map whose response depends on the level of the latent outcome is outside it.
And $m_c(\ell) = 0$ for $\ell < 0$ is no anticipation in the recording layer:
the recording system changes when treatment arrives, not before. Unit by category
terms $\mu_{ic}$ are unrestricted, so any time invariant difference in recording
practice across units is absorbed and only the treatment-induced change matters.

\subsection{Pre trend tests have power equal to size}

Let $\Delta_c(g,t)$ denote the deviation from parallel trends in the latent
layer. The following is a coupling argument: it exhibits, for any recording
response, a structure with no recording response whatever that generates the
same observed data.

\begin{lemma}[Observational equivalence]
\label{lem:oe}
Fix a structure satisfying Assumption \ref{as:rec} with latent parallel trends
and recording paths $m_c$. The structure with no recording response and latent
deviation
\begin{equation}
\label{eq:equivdev}
\Delta_c(g,t) = m_c(t-g)\,\mathbf 1\{t \ge g\}
\end{equation}
generates the same law of the observed panel, cell by cell. Its pre-treatment deviations are exactly zero.
\end{lemma}

\begin{corollary}[No pre-treatment diagnostic separates the two]
\label{cor:blind}
Let $\varphi$ be any test that is a measurable function of the observed panel
and the cohort labels. Its rejection probability is the same under both
structures. In particular a pre-trend test, a placebo in time test and a
leads and lags plot reject a recording response with probability no greater than
their nominal level, and exactly equal to it when the test is similar on the
pre-treatment null.
\end{corollary}

The online Appendix answers four objections to Corollary \ref{cor:blind}: that
it restates a familiar fact, that it applies to the test this paper goes on to
propose, that pre-treatment parallelism has long been known not to guarantee its
post-treatment counterpart, and how the construction differs from the one it
resembles.

\subsection{Sensitivity analysis calibrated from pre trends}

\begin{proposition}[The relative magnitudes family is empty against a recording response]
\label{prop:rm}
Under Assumption \ref{as:rec} with $m_c \not\equiv 0$, the population identified
set for the latent effect under $\Delta^{\mathrm{RM}}(\bar M)$ is the singleton
$\{\vartheta_c(\ell) + m_c(\ell)\}$ for every $\bar M \ge 0$, and it excludes the
latent effect by exactly the recording wedge.
\end{proposition}

Proposition \ref{prop:rm} says the relative magnitudes class is empty against a
recording response: the class scales a pre-period quantity that Corollary
\ref{cor:blind} makes identically zero, so it admits no deviation at any value
of $\bar M$, and the smoothness class calibrated from the pre-period inherits
the same failure through its estimated bound. The online Appendix develops both
classes and the calibration the reference software uses.

What is prior here should be marked. That a smoothness constant cannot be read
off the data while honest coverage is retained is the nonadaptivity result of
\citet{low1997nonparametric} and \citet{cailow2004adaptation}, stated for
econometric practice by \citet{armstrong2020simple} and restated in a
difference-in-differences framing by \citet{freeman2026local}, and
\citet{arora2026trading} note in prose that uninformative pre trends leave the
relative magnitudes class with intervals that fail to widen. Proposition
\ref{prop:rm} improves on none of that. What it adds is the mechanism: under a
recording response the quantity being scaled is zero by construction rather
than small by chance, so the collapse is exact, holds at every $\bar M$, and is
a property of the structure rather than of the realized sample.

\subsection{What coverage does}

Coverage of the latent effect by a pre-period calibrated honest confidence set
falls from $0.955$ at $N = 400$ to $0.000$ from $N = 25{,}600$ on, while a bound
set from the recording institution holds at $1.000$ at every sample size and the
recorded effect is covered throughout. The calibration input falls at the
sampling rate while the wedge does not move, so every relative magnitudes
interval eventually excludes the latent effect and larger $\bar M$ only delays
the crossing. More data makes it worse.
The online Appendix gives the full table, the three blocks it contains and the
reference implementation the sets are computed with.

\subsection{Scope}

Three limits bound the claims above. Assumption \ref{as:rec} restricts the
recording map to enter additively with a path indexed by event time; a map that
responds to the level of the latent outcome is not covered. Lemma \ref{lem:oe}
and Corollary \ref{cor:blind} do not require latent parallel trends, but
Proposition \ref{prop:rm} is stated for a single adoption cohort, and the
extension to staggered designs is not proved here, though the online Appendix
reports the same collapse in a three cohort design. And the argument is about
what pre-treatment data can reveal; it says nothing against sensitivity analysis
whose bound comes from outside the event-study coefficients, which is the
practice the coverage experiment of the online Appendix supports and which
\citet{armstrong2020simple} already recommend.

The last limit is where the result stops rather than where it fails, and it is
the one most likely to be mistaken for a refutation. Corollary \ref{cor:blind}
concerns diagnostics computed from the outcome panel and the cohort labels. It
says nothing about identifying a recording response from data on the recording
process itself, and applied work already does this successfully:
\citet{ody2019readmissions} separate the Hospital Readmissions Reduction Program
from a coincident change in electronic transaction standards using the
distribution of diagnosis counts per claim rather than the readmission series.
Auxiliary information about the record is not a counterexample to the corollary.
It is what the corollary says is required.

What survives is a design problem. If the contamination is a recording response,
the pre-period is silent, and information has to come from somewhere else. The
remainder of this paper takes it from the cross section of outcome categories,
which differ in how much discretion the recording system has over them, and
addresses the inference problem that arises when the number of clusters is
small.

\section{Setup and the null}
\label{sec:setup}

Units $i = 1,\dots,N$ are grouped into clusters $g = 1,\dots,G$; periods are
$t = 1,\dots,T$; outcome categories are $c = 1,\dots,C$ with known exposure
indices $d_1,\dots,d_C$. Treatment is absorbing with cohort $G_i \in
\{1,\dots,T\} \cup \{\infty\}$ and $D_{it} = \mathbf{1}\{t \ge G_i\}$.

\begin{assumption}[category balanced design]
\label{as:B}
There is a common cell set $\mathcal{C} \subseteq \{1,\dots,N\} \times
\{1,\dots,T\}$ with $|\mathcal{C}| = n_0$ on which every category is observed,
and the treatment path $D_{it}$, the cohort $G_i$ and the cluster map $i \mapsto
g(i)$ do not depend on $c$.
\end{assumption}

Assumption \ref{as:B} rules out designs in which categories are measured on
different subsamples or at different frequencies. It is verifiable from the data
layout. It is a hypothesis of Lemma \ref{lem:sep} rather than of Theorem
\ref{thm:exact}: when it fails, $\hat{\gamma}$ in \eqref{eq:pooled} is no longer
the fixed linear combination of per category event studies that Lemma
\ref{lem:sep} gives and each category has to be fitted on its own cells, while
the randomization of Section \ref{sec:perm} acts in the category direction and
needs only Assumption \ref{as:CE}. Section \ref{sec:app} reports one application in which
Assumption \ref{as:B} fails through a cell pattern that depends on the category
label, which breaks Assumption \ref{as:CE} as well.\footnote{Fitting each category on its own
cells, the label permutation test of \eqref{eq:null} rejects at $0.047$ when
the missing cells are tied to the category label and at $0.049$ when they are
independent of it, against a nominal $0.05$ over $1{,}200$ replications with a
Monte Carlo standard error of $0.006$
(\texttt{verification\_extra6\_3.json} in \texttt{results/inference/}). With the
observation rate falling from $1.00$ to $0.40$ across categories ordered by
$d_c$, the rate is $0.040$, against $0.053$ for a control that assigns the same
five rates to a category drawn at random in each replication, over $1{,}500$
replications with a Monte Carlo standard error of $0.006$
(\texttt{verification\_extra7\_8.json}, same directory). The second
departure is downward, so the test stays at or below its nominal level and the
cost falls on power.}

Let $L$ collect the included event time horizons with $|L| = q$, omitting
$\ell = -1$ and binning endpoints, let $X \in \mathbb{R}^{n_0 \times q}$ hold the
corresponding indicators, and let $P \subseteq \{1,\dots,q\}$ with $|P| = m$ be
the post-treatment block. With $y_{itc}$ the outcome and $\delta_c = d_c -
\bar{d}$ the centered exposure index, the pooled gradient regression is
\begin{equation}
y_{itc} = a_{ic} + p_{tc} + \sum_{\ell \in L} \beta_\ell X_{(it),\ell}
  + \sum_{\ell \in L} \gamma_\ell \delta_c X_{(it),\ell} + u_{itc},
\label{eq:pooled}
\end{equation}
estimated by OLS with the CR1 cluster-robust covariance clustered on $g$. The
hypothesis under test throughout is the joint null
\begin{equation}
H_0: \gamma_P = 0,
\label{eq:null}
\end{equation}
that the post-treatment coefficients do not vary with exposure. This is a weak
null on a parameter vector, not a sharp null of no effect for every unit, and
the distinction matters for the randomization arguments in Sections
\ref{sec:perm} and \ref{sec:cluster}: we state separately for each procedure
which invariance delivers validity and under what condition.

The cross-category structure enters \eqref{eq:pooled} as a single interaction
contrast. The procedure tests one joint hypothesis, not $C$ separate hypotheses,
so no multiplicity adjustment arises and the error rate controlled is the
per comparison rate for \eqref{eq:null}.

Write $\Pi$ for the two-way within projection on $\mathcal{C}$,
$\tilde{X} = \Pi X$, $S = \tilde{X}'\tilde{X}$, and $b_c = S^{-1}\tilde{X}'\Pi y_c$ for the
within category event-study coefficient vector.

\begin{lemma}[separability]
\label{lem:sep}
Under Assumption \ref{as:B} and $\operatorname{rank}(\tilde{X}) = q$, the OLS
estimators in \eqref{eq:pooled} satisfy $\hat{\beta} = C^{-1}\sum_c b_c$ and
$\hat{\gamma} = \sum_c \delta_c b_c / \sum_c \delta_c^2$.
\end{lemma}

The gradient estimator is a fixed linear combination of the per category event
studies, which is what makes the category level randomization of Section
\ref{sec:perm} available.

The $b_c$ are within category event studies estimated by OLS, which under
staggered adoption with heterogeneous effects need not recover a convex
combination of cohort specific effects
\citep{dechaisemartin2020twoway,sunabraham2021}. Replacing each $b_c$ by a
heterogeneity robust estimator \citep{callaway2021did,borusyak2024revisiting},
or by another of the alternatives surveyed in \citet{roth2023trending}, leaves
the randomization arguments of Sections \ref{sec:perm} and \ref{sec:cluster}
intact, because those require only that the category level statistic be
order invariant in the sense of Theorem \ref{thm:exact}. The exact null law of
Section \ref{sec:anatomy} is specific to the CR1 Wald statistic from
\eqref{eq:pooled} and would have to be rederived for another estimator.

Failing to reject \eqref{eq:null} does not establish that the identifying
assumption holds. Section \ref{sec:power} gives the minimum detectable response
for each design, and the applications in Section \ref{sec:app} report it
alongside every nonrejection.

\section{The exact null law of the cluster robust Wald statistic}
\label{sec:anatomy}

The size distortion has a closed form. We isolate it under an idealization that
is used only for this diagnosis and is not maintained elsewhere in the paper.

Let $\tilde{X}_g$ collect the rows of $\tilde{X}$ belonging to cluster $g$, let
$u_{gc}$ be the corresponding errors in \eqref{eq:pooled}, let $e_{gc} =
\tilde{y}_{gc} - \tilde{X}_g b_c$ be the within category residuals, and let
$n = C n_0$ be the number of observations in \eqref{eq:pooled}. The cluster
score of the gradient block and its fitted counterpart are
\begin{equation}
a_g = \frac{\tilde{X}_g'\sum_c \delta_c u_{gc}}{\sum_c \delta_c^2},
\qquad
a^f_g = \frac{\tilde{X}_g'\sum_c \delta_c e_{gc}}{\sum_c \delta_c^2},
\qquad
\eta_g = S^{-1}a_g, \qquad \eta^f_g = S^{-1}a^f_g ,
\label{eq:eta}
\end{equation}
all in $\mathbb{R}^q$. Lemma \ref{lem:sep} gives $\hat{\gamma} - \gamma =
\sum_g \eta_g$, and the CR1 covariance of $\hat{\gamma}$ is $\omega \sum_g
\eta^f_g \eta^{f\prime}_g$ with $\omega = \{G/(G-1)\}\{(n-1)/(n-2q)\}$. A
subscript $P$ denotes the subvector or submatrix indexed by the post block, and
$\bar{\eta}_P = G^{-1}\sum_g \eta_{gP}$.

\begin{assumption}[ideal cluster experiment]
\label{as:I}
(i) Under $H_0$ in \eqref{eq:null}, the cluster score contributions
$\eta_1,\dots,\eta_G$ are i.i.d.\ $N(0,\Sigma)$ with $\Sigma_{PP} \succ 0$.
(ii) Every cluster carries the same design weight, $S_g = S/G$ for all $g$,
where $S_g = \tilde{X}_g'\tilde{X}_g$.
\end{assumption}

Part (i) holds when the errors in \eqref{eq:pooled} are Gaussian, independent
across clusters, and identically distributed across clusters given the design.
Part (ii) fails in a staggered design with a large never-treated share;
Definition \ref{def:geff} handles that case.

\begin{proposition}[exact null law]
\label{prop:anatomy}
Let Assumptions \ref{as:B} and \ref{as:I} hold, let
$\operatorname{rank}(\tilde{X}) = q$, let $H_0$ in \eqref{eq:null} hold, and let
$G - m \ge 1$. Put $\omega_2 = (n-2q)/(n-1) = G/\{(G-1)\omega\}$, let $\hat{S}_P$
be the sample covariance of the $\eta_{gP}$ and $T^2 = G \bar{\eta}_P'
\hat{S}_P^{-1} \bar{\eta}_P$. Then the CR1 Wald statistic $W$ for
\eqref{eq:null} satisfies $W = \omega_2 T^2$ identically, and
\begin{equation}
\frac{(G-m)\,W}{\omega_2 (G-1) m} \sim F(m, G-m)
\label{eq:Fexact}
\end{equation}
exactly for every $G$ and $m$ with $G - m \ge 1$. If instead $G - 1 < m$, the
CR1 meat $\sum_g \eta^f_{gP}\eta^{f\prime}_{gP}$ has rank at most $G-1$, the
Moore--Penrose version of $W$ has no Hotelling representation, and
\eqref{eq:Fexact} is void.
\end{proposition}

The chi-square calibration treats the estimated meat matrix as known. Under
Assumption \ref{as:I} the cost of that substitution is exactly the $F$
denominator in \eqref{eq:Fexact}, which has $G - m$ degrees of freedom. The
statement is a degrees of freedom result: nothing about the errors beyond
Assumption \ref{as:I}(i) enters, and $\omega_2$ lies between $0.995$ and $0.9997$ across the designs we
study and equals $0.9989$ at the headline
design of Section \ref{sec:sim} ($G = 51$, $C = 5$). The factor $(n-2q)/(n-1)$
counts the $2q$ interacted event time regressors in \eqref{eq:pooled} and not
the $C(N+T-1)$ absorbed unit by category and period by category effects, so a
user calibrating through \eqref{eq:Fexact} must supply the $\omega_2$ that
matches the degrees of freedom convention of their own CR1 implementation.

Inverting \eqref{eq:Fexact} for a critical value gives a reference distribution
that is exact for the null \eqref{eq:null} under Assumption \ref{as:I}, whose
part (ii) requires every cluster to carry the same design weight. Where that
part fails the calibration is neither exact nor conservative. In the staggered
design of Section \ref{sec:sim} with $C = 5$, referring $W$ to \eqref{eq:Fexact}
at the true $G$ gives null rejection rates of $0.174$, $0.160$, $0.124$ and
$0.092$ at $G = 20$, $30$, $51$ and $100$ against the nominal $0.05$, over
$2{,}000$ replications per cell (\texttt{anatomy\_table.csv}, which is also the
source for $\omega_2$ above). The version evaluated at the $G_{\mathrm{eff}}$ of
Definition \ref{def:geff} is an approximation whose accuracy Section
\ref{sec:sim} measures.
The $F$ family in \eqref{eq:Fexact} is the one that \citet{tipton2015small} and
\citet{pustejovsky2018small} reach by matching denominator degrees of freedom to
the data; under Assumption \ref{as:I} no matching is needed, because
\eqref{eq:Fexact} holds exactly with $G - m$.

\begin{corollary}[when the chi-square calibration is valid]
\label{cor:asym}
Let $\mathrm{size}(G,m)$ denote the null rejection probability of the nominal
level-$\alpha$ chi-square test under the conditions of Proposition
\ref{prop:anatomy}. Assume $G(1-\omega_2) = O(1)$ in part (i) and
$\sqrt{m}\,|\log \omega_2| \to 0$ in part (ii), where $1 - \omega_2 =
(2q-1)/(n-1)$.

\noindent (i) For fixed $m$, $\mathrm{size}(G,m) = \alpha + c_m f_m(c_m)
\big[(m + c_m)/\{2(G-1)\} - (1-\omega_2)\big] + O(G^{-2})$ where $c_m =
\chi^2_{m,1-\alpha}$ and $f_m$ is the $\chi^2_m$ density, so
$\mathrm{size}(G,m) \to \alpha$ as $G \to \infty$. The correction is positive,
and the chi-square test over rejects at all large $G$, when
$2(G-1)(1-\omega_2) < m + c_m$.

\noindent (ii) Let $m \to \infty$ with $m/G \to 0$ and $\chi_{m,G} =
\sqrt{m/2}\,\log\{(G-1)/(G-m)\} \to \chi \in [0,\infty]$. Then
$\mathrm{size}(G,m) \to \Phi(\chi - z_{1-\alpha})$, and $\chi_{m,G} =
m^{3/2}/(\sqrt{2}\,G) + o(1)$ whenever $m^{5/2}/G^2 \to 0$. If instead
$m/G \to r \in (0,1)$, then $\mathrm{size}(G,m) \to 1$.
\end{corollary}

Part (ii) is the operative statement. The chi-square calibration is
asymptotically valid if and only if $m^{3/2}/G \to 0$, that is $m = o(G^{2/3})$. The online Appendix plots the exact size implied by \eqref{eq:Fexact} against the number of restrictions at five cluster counts, with the measured rates of Table \ref{tab:size} overlaid.

The limit in Corollary \ref{cor:asym}(ii) is not accurate at the $m$ this paper
works with: along the sequence with $\chi = 0.5$, the exact size from \eqref{eq:Fexact}
exceeds $\Phi(\chi - z_{1-\alpha})$ by $5.95$ percentage points at $m = 13$,
$4.21$ at $m = 25$ and $2.91$ at $m = 50$, an error falling like $m^{-1/2}$
(\texttt{anatomy\_asymptotics.csv}). Every admissible-$m$ figure reported below
is computed from \eqref{eq:Fexact} and not from the limit.

\begin{definition}[effective clusters]
\label{def:geff}
Let $J_P$ be the $q \times m$ selection matrix whose columns are the standard
basis vectors indexed by $P$, so that $[S^{-1}]_{PP} = J_P'S^{-1}J_P$. For a
contrast $w \in \mathbb{R}^m$ put $a = S^{-1}J_P w$ and $\varphi_g = a'S_g a
\ge 0$. Then $G_{\mathrm{eff}}(w) = (\sum_g \varphi_g)^2 / \sum_g
\varphi_g^2$.
\end{definition}

The online Appendix derives $G_{\mathrm{eff}}$ under homoskedastic errors
independent across clusters and reports the value every design here takes.

\section{The label permutation test}
\label{sec:perm}

Write $\mathcal{Y}_c \in \mathbb{R}^{n_0}$ for the category-$c$ data array,
$\mathcal{D} = (\mathcal{Y}_1,\dots,\mathcal{Y}_C)$, and $\mathcal{Z}$ for the
design, which is common to all categories under Assumption \ref{as:B}. Let
$\mathcal{T}(\mathcal{D},\delta)$ be a statistic with large values evidence
against $H_0$ in \eqref{eq:null}.

The statistic is computed on loading normalized arrays. Let
$\hat{\Sigma}_{ck}$ be the sample covariance of the two-way demeaned categories
$c$ and $k$ over the untreated cells $\{(i,t) \in \mathcal{C}: D_{it} = 0\}$,
and put $L_{ck} = \log|\hat{\Sigma}_{ck}|$. The log triad loading estimate is
$\hat{\rho}_c \propto \exp(\hat{a}_c)$, scaled to geometric mean one, with
\begin{equation}
\hat{a}_c = \frac{1}{C-1}\sum_{k \ne c} L_{ck}
 \;-\; \frac{1}{2(C-1)(C-2)} \sum_{k \ne l,\; k,l \ne c} L_{kl} ,
\qquad C \ge 3 .
\label{eq:triad}
\end{equation}
Under a one-factor structure with $\hat{\Sigma}_{ck} \to_p \theta_c\theta_k s^2$
for $c \ne k$, the limit of $L_{ck}$ is $\log|\theta_c| + \log|\theta_k| +
\log s^2$, and the two averages in \eqref{eq:triad} isolate $\log|\theta_c|$ up
to a constant common to all $c$, so $\hat{\rho}_c \to_p |\theta_c| /
(\prod_j |\theta_j|)^{1/C}$. There is no reference category and no ordering
convention in \eqref{eq:triad}, which is what Lemma \ref{lem:equi} records.

\begin{definition}[label permutation test]
\label{def:perm}
Normalize each array, $\mathcal{Y}^{\dagger}_c = \mathcal{Y}_c/\hat{\rho}_c$,
and drop every category with $|\hat{\rho}_c| \le \rho_{\min}$, where
$\rho_{\min} = 0.1$ throughout; write $C_0$ for the count of Section
\ref{sec:setup}, $C$ for the number of categories that survive the screen, and
$\mathcal{D}^{\dagger}$ for the surviving normalized arrays. For
$\pi \in S_C$ let $\delta_\pi = (\delta_{\pi(1)},\dots,\delta_{\pi(C)})$ and
$T(\pi) = \mathcal{T}(\mathcal{D}^{\dagger},\delta_\pi)$. With reference set
$\mathcal{R} = S_C$ and $M = C!$, the test rejects at level $\alpha$ when
$\hat{p} = M^{-1}\#\{\pi \in \mathcal{R}: T(\pi) \ge T(\mathrm{id})\} \le
\alpha$. When $C! > B$ the reference set is instead $\mathcal{R} =
\{\mathrm{id},\pi_2,\dots,\pi_B\}$, with $\pi_2,\dots,\pi_B$ drawn
independently and uniformly from $S_C$, and $M = B$.
\end{definition}

The randomization is over the assignments of the exposure labels to the
normalized category data arrays, holding the arrays and the design fixed.
Nothing is permuted at the unit, cluster or period level. Every run reported
below sets $B = 999$, so the reference set is all of $S_C$ at $C \le 6$ and a
Monte Carlo sample at $C = 8$. Enumeration costs $C!$ evaluations of
$\mathcal{T}$ and is the binding constraint well before the statistic is: $720$
at $C = 6$, $5{,}040$ at $C = 7$ and $40{,}320$ at $C = 8$, against the $999$ a
Monte Carlo reference set uses at any $C$. Sampling the group costs resolution
rather than validity. Theorem \ref{thm:exact} holds on either reference set,
the attainable floor becomes $1/B$ by Proposition \ref{prop:resolution}(i), and
the Monte Carlo $p$ value has standard error at most $1/(2\sqrt{B})$, which is
$0.016$ at $B = 999$. The batched implementation evaluates all $B$ label vectors
in one pass over the cluster scores, so the cost is linear in $B$ and does not
grow with $C!$.

\begin{lemma}[equivariance of the normalization]
\label{lem:equi}
Under Assumption \ref{as:B}, relabelling the arrays relabels the loading
estimates identically: writing $\mathcal{D}_\sigma =
(\mathcal{Y}_{\sigma(1)},\dots,\mathcal{Y}_{\sigma(C)})$, one has
$\hat{\rho}_c(\mathcal{D}_\sigma) = \hat{\rho}_{\sigma(c)}(\mathcal{D})$ for
every $\sigma \in S_C$. The retained set $\{c: |\hat{\rho}_c| > \rho_{\min}\}$
is relabeled with them, so $\mathcal{D} \mapsto \mathcal{D}^{\dagger}$ is
$S_C$-equivariant.
\end{lemma}

A statistic that is order invariant in the normalized arrays is therefore
order invariant in the raw arrays, and Theorem \ref{thm:exact} below applies to
$\mathcal{D}^{\dagger}$ unchanged.

The online Appendix scopes \eqref{eq:triad}: two conditions restricting the
procedure rather than the panel, and the one-factor structure it inverts,
which the crime designs support and the health design does not.

\begin{assumption}[category exchangeability]
\label{as:CE}
Conditional on $\mathcal{Z}$,
\[
(\mathcal{Y}_{\pi(1)},\dots,\mathcal{Y}_{\pi(C)})
  \overset{d}{=} (\mathcal{Y}_1,\dots,\mathcal{Y}_C)
\quad\text{for every } \pi \in S_C,
\]
and $d_1,\dots,d_C$ are fixed constants.
\end{assumption}

Assumption \ref{as:CE} rules out categories with different error variances and
any dependence of the labels on the data. It restricts the category direction
only: heteroskedasticity across units and periods, serial correlation, and
within cluster dependence of any form are unrestricted, and $G$ is unrestricted.
Normalizing the outcome does not restore it under fixed heterogeneous loadings.
Dividing category $c$ by its true loading $\theta_c$ equalizes the factor
component and leaves $\operatorname{Var}(\varepsilon_c/\theta_c) =
\sigma^2/\theta_c^2$ category specific, and at $C = 5$ with
$\theta = (0.4,0.7,1.0,1.3,1.6)$ the label permutation test of Definition
\ref{def:perm} run on the true loading normalization has size $0.030$ against
the nominal $0.05$, with Monte Carlo standard error $0.006$
(\texttt{verification.json}, entry \texttt{V24}). Exchangeability
does hold when the loadings are exchangeable draws independent of $d$, which is
what the simulation of Section \ref{sec:sim} imposes, and by Lemma
\ref{lem:equi} estimating them through \eqref{eq:triad} preserves it.
Assumption \ref{as:CE} is not implied by $H_0$ in \eqref{eq:null} and it is not
directly testable, though its leading failure mode has an observable signature:
Section \ref{sec:sim} prices the size distortion as a function of the ratio of
the largest to smallest category error variance.

\begin{theorem}[exact finite-sample validity]
\label{thm:exact}
Suppose Assumption \ref{as:CE} holds and $\mathcal{T}$ is order invariant, in
the sense that it is a function of the multiset
$\{(\delta_c,\mathcal{Y}^{\dagger}_c)\}_{c=1}^C$ and not of the order in which
the pairs are listed. Then for every $N$, $T$, $G$ and $C$, the test of
Definition \ref{def:perm} for the null \eqref{eq:null} under the randomization
over its reference set $\mathcal{R}$ satisfies $\Pr(\text{reject}) \le \alpha$.
If the distinct values among $\{T(\pi): \pi \in \mathcal{R}\}$ are almost surely
attained with a common multiplicity $\mathfrak{m}$, then $\Pr(\text{reject}) =
\mathfrak{m} \lfloor \alpha M/\mathfrak{m} \rfloor / M$. Here $\mathfrak{m} = 2$ for the
two sided Wald statistic with $\{\delta_c\}$ symmetric about zero, by
Proposition \ref{prop:resolution}(ii), and $\mathfrak{m} = 1$ for the one-sided directed statistic. At $C = 5$ and $\alpha = 0.05$ the two give $0.05$ alike.
\end{theorem}

Exactness here is finite-sample exactness for the null \eqref{eq:null} under the
randomization distribution of Definition \ref{def:perm}, and it requires
Assumption \ref{as:CE}. No condition on $G$, on the error distribution, or on
the within cluster dependence structure is used. Order invariance holds for the
OLS estimates and CR1 covariance in \eqref{eq:pooled}, and by Lemma
\ref{lem:equi} it survives the normalization of Definition \ref{def:perm}. It
fails when the normalization is not $S_C$-equivariant, which a reference category
rule keyed to $\arg\min_c d_c$ is not. A position keyed rule that acts as an
order invariant rescaling times a scalar common to all categories remains valid,
because $\hat{\gamma}$ and its CR1 covariance in \eqref{eq:pooled} are
homogeneous of degree one and two in the outcome and the Wald statistic is
therefore unchanged by a common rescaling: a normalization whose reference
column is chosen by the position of the largest permuted label has size $0.052$
against the nominal $0.05$, with Monte Carlo standard error $0.005$
(\texttt{verification.json}).

A randomized version has size exactly $\alpha$. It rejects with probability one
when $\hat{p} \le \alpha_-$ and with probability $(\alpha-\alpha_-)/
(\alpha_+-\alpha_-)$ when $\hat{p} = \alpha_+$, where $\alpha_-$ is the largest
attainable value of $\hat{p}$ not exceeding $\alpha$ and $\alpha_+$ the smallest
value exceeding it. No procedure reported below uses it.

\begin{proposition}[attainable $p$ values]
\label{prop:resolution}
(i) On the exhaustive reference set, $\hat{p} \ge |\mathrm{Stab}(\delta)|/C!$,
where $\mathrm{Stab}(\delta) = \{\pi \in S_C : \delta_\pi = \delta\}$ has order
the product of the factorials of the sizes of the tied groups of
$d_1,\dots,d_C$, so the bound is $1/C!$ when $d$ has no ties. On the Monte Carlo
reference set of Definition \ref{def:perm}, $\hat{p} \ge 1/B$. (ii) If
$\mathcal{T}$ is even in the labels, so that
$\mathcal{T}(\mathcal{D}^{\dagger},-\delta) =
\mathcal{T}(\mathcal{D}^{\dagger},\delta)$, and the multiset $\{\delta_c\}$ is
symmetric about zero and not identically zero, then $\hat{p} \ge
2|\mathrm{Stab}(\delta)|/C!$ on the exhaustive reference set. Equally spaced $d$
is the leading sufficient case for the symmetry. (iii) On the exhaustive
reference set with $d$ untied, a level-$\alpha$ label permutation test therefore
requires $C! \ge 2/\alpha$ for a statistic even in the labels with
$\{\delta_c\}$ symmetric about zero, that is $C \ge 5$ at $\alpha = 0.05$, and
$C! \ge 1/\alpha$ otherwise, that is $C \ge 4$.
\end{proposition}

The Wald statistic is even in the labels because $\hat{\gamma}(-\delta) =
-\hat{\gamma}(\delta)$ and its covariance estimator is quadratic in the cluster
scores. Proposition \ref{prop:resolution} is a property of the randomization
group, not of any implementation. With $C = 3$ the two sided Wald has three
attainable $p$ values, the smallest being $1/3$, and the test can never reject
at conventional levels. Section \ref{sec:sim} reports the resulting rejection
rate of exactly $0.000$.

\section{Cluster level randomization and the bootstrap}
\label{sec:cluster}

Two procedures remain available when $C \le 4$. Both randomize over clusters
rather than category labels, so Proposition \ref{prop:resolution} does not bind.

The first randomizes sign changes over $\mathcal{E} = \{-1,+1\}^G$. Let
$\gamma^r$ be the restricted estimate, with $\gamma^r_P = 0$ and
$\gamma^r_{\mathrm{pre}} = (S_{\mathrm{pre},\mathrm{pre}})^{-1}
(S\hat{\gamma})_{\mathrm{pre}}$, let $z^r_g = \sum_c \delta_c \tilde{X}_g'
\{\tilde{y}_{gc} - \tilde{X}_g(\hat{\beta} + \delta_c \gamma^r)\}$ be the
cluster-$g$ score of the gradient block at that restricted fit, and put
\begin{equation}
\zeta_g = \frac{(S^{-1}z^r_g)_P}{\sum_c \delta_c^2} \in \mathbb{R}^m ,
\qquad \sum_g \zeta_g = \hat{\gamma}_P .
\label{eq:zeta}
\end{equation}
The test randomizes $\mathcal{Q}(\epsilon) = (\sum_g \epsilon_g \zeta_g)'
\hat{\Omega}^{-1}(\sum_g \epsilon_g \zeta_g)$ over $\epsilon \in \mathcal{E}$,
with $\hat{\Omega}$ proportional to $\sum_g \zeta_g\zeta_g'$, and rejects when
$\mathcal{Q}(\mathbf{1}) = \hat{\gamma}_P'\hat{\Omega}^{-1}\hat{\gamma}_P$ falls
in the upper $\alpha$ tail of $\{\mathcal{Q}(\epsilon)\}$. The CR1 scale factor
in $\hat{\Omega}$ is common to every $\epsilon$ and cancels from the comparison,
and $\hat{\Omega}$ itself does not depend on $\epsilon$ because
$\epsilon_g^2 = 1$. Write $\zeta^0_g$ for the same functional of the true errors
under \eqref{eq:null}.

\begin{theorem}[cluster sign-change test]
\label{thm:signflip}
If, under $H_0$ in \eqref{eq:null}, the cluster contributions are jointly
sign symmetric, meaning $(\epsilon_1 \zeta_1,\dots,\epsilon_G \zeta_G)
\overset{d}{=} (\zeta_1,\dots,\zeta_G)$ for every $\epsilon \in \mathcal{E}$,
then the sign-change test is exactly level $\alpha$ in finite samples, and so is
its Monte Carlo version using $B-1$ random sign vectors together with the
identity. Sign symmetry holds for the infeasible statistic built from the true
errors under \eqref{eq:null} whenever the cluster error contributions are
independent across $g$ and symmetric about zero. Under Assumption
\ref{as:I}(ii) the feasible and infeasible contributions coincide exactly on
$P$, so the feasible test is then exactly level $\alpha$ in finite samples. The feasible statistic differs
from the infeasible one by $O_p(G^{-1/2})$ and is asymptotically level $\alpha$
as $G \to \infty$ under conditions (W1) to (W3) stated in the online Appendix.
\end{theorem}

Under Assumption \ref{as:I}(ii) the gap between the two statistics is zero,
because $\zeta_g - \zeta^0_g = -(S^{-1}S_g\gamma^r)_P = -\gamma^r_P/G = 0$ by
construction of $\gamma^r$, and the measured gap is exactly zero at $G = 12$,
$24$, $48$, $96$ and $192$ (\texttt{verification.json}). Outside that condition the finite-sample statement
applies to the infeasible statistic alone, and the feasible statistic inherits
validity only asymptotically in $G$. The argument for that limit in the
online Appendix is a conditional bootstrap central limit theorem rather than an
approximate symmetry argument, because the sign-change group has cardinality
$2^G$. Section \ref{sec:sim} measures the finite-sample gap.

The statistic $\mathcal{Q}$ is even in $\epsilon$, so its $2^G$ values pair and
the smallest attainable $p$ value is $2^{1-G}$ rather than $2^{-G}$. An
exhaustive enumeration at $G = 7$ gives $64$ distinct values out of $128$ and a
minimum $p$ value of $2/128$ (\texttt{verification.json}, entry \texttt{E14}).
That floor does not bind at any $G$ reported below, the smallest of which is
$20$.

The second procedure is the restricted wild cluster bootstrap of
\citet{cameron2008bootstrap}, drawing Rademacher weights at the cluster level
with \eqref{eq:null} imposed on the residuals. The implementation partials the
unit by category and period by category effects out once and applies the weights
to the cluster residuals of the partialled out regression, rather than
reabsorbing those effects on every draw; the two versions of the statistic
differ by $O_p(G^{-1})$. No finite-sample result is claimed for this procedure.
\citet{djogbenou2019asymptotic} give conditions for asymptotic validity of the
wild cluster bootstrap in a linear regression with clustered errors:
independence across clusters, moment conditions on the cluster level scores, and
a number of clusters growing without bound with no cluster dominating. Section
\ref{sec:sim} reports its size in the designs run there and claims nothing
beyond that.

\citet{mackinnon2017wild} report simulations in which the cluster robust
$t$-test over rejects more when the $G = 50$ cluster sizes are proportional to
United States state populations than when they are balanced, in which the wild
cluster bootstrap corrects most of that over rejection, and in which it fails
when only a few clusters are treated. Their designs vary the cluster size
distribution and the number of treated clusters, not the ratio $m/G$ that drives
Corollary \ref{cor:asym}. Every cluster in the simulation of Section
\ref{sec:sim} holds one unit observed for $T = 120$ periods, so that simulation
cannot test their failure mode either.

The online Appendix sets out the six step recipe for running the three
procedures, and names the file in the replication archive that implements it.

\section{Power, the minimum detectable response, and the optimal statistic}
\label{sec:power}

Let the alternative be $\gamma_\ell = \kappa\, k(\ell)$ for a response path
$k(\cdot) \ge 0$ and a scalar magnitude $\kappa > 0$, so that $H_1: \kappa > 0$.
Under local alternatives the noncentrality of the Wald statistic depends on the
design through two scalars: $\Sigma_\delta^2 = \sum_c (d_c - \bar{d})^2$, a
property of the exposure index, and the timing quantity
\begin{equation}
I_k = k_P'\big([S^{-1}]_{PP}\big)^{-1}k_P ,
\label{eq:ik}
\end{equation}
where $k_P = (k(\ell))_{\ell \in P}$ is the response shape on the
post-treatment block. In the single cohort design with $N$ units, an adopting
share $p$, $T_0$ omitted reference periods and the constant path
$k \equiv \mathbf{1}$, $[S^{-1}]_{PP} = \{Np(1-p)\}^{-1}(I_m +
\mathbf{1}\mathbf{1}'/T_0)$ and $I_k = Np(1-p)\,mT_0/(m+T_0)$, which saturates
in $T$: with a single omitted period the information about a level shift never
exceeds $Np(1-p)$, however long the panel. The staggered design of Section
\ref{sec:sim} gives $I_k = 23.62$.

\begin{corollary}[closed-form minimum detectable response]
\label{cor:mde}
Let $\lambda^*(\alpha,\beta,m)$ solve $\Pr\{\chi^2_m(\lambda^*) >
\chi^2_{m,1-\alpha}\} = 1-\beta$. With a linear wedge and homoskedastic
categories,
\begin{equation}
\kappa^{\mathrm{Wald}}_{\mathrm{MDE}} = \sigma
\sqrt{\frac{\lambda^*(\alpha,\beta,m)}{\Sigma_\delta^2\, I_k}},
\qquad
\kappa^{\mathrm{dir}}_{\mathrm{MDE}} = \sigma\,
\frac{z_{1-\alpha}+z_\beta}{\sqrt{\Sigma_\delta^2\, I_k}} .
\label{eq:mde}
\end{equation}
\end{corollary}

Both $\Sigma_\delta^2$ and $I_k$ are computable before any outcome is observed.
The third input, $\sigma^2$, is the $\delta^2$-weighted mean $\sum_c \delta_c^2
\sigma_c^2/\sum_c \delta_c^2$ of the category residual variances $\sigma_c^2$ of
the loading normalized outcome, and it is the only quantity requiring pilot
data; a researcher estimates it by the residual variance of the within category
regression in \eqref{eq:pooled} run on the normalized outcome over untreated
cells. Every figure below sets $\beta = 0.20$, so power is $1-\beta = 0.80$, and
$\lambda^*(0.05,0.20,13) = 17.85$. The expression for
$\kappa^{\mathrm{Wald}}_{\mathrm{MDE}}$ is stated against the chi-square critical value, which Proposition \ref{prop:anatomy} shows is the wrong one; it
therefore understates the minimum detectable response of any correctly
calibrated Wald test. Replacing $\chi^2_{m,1-\alpha}$ by the exact critical
value from \eqref{eq:Fexact} raises the required noncentrality from $17.85$ to
$41.46$ at $G_{\mathrm{eff}} = 36.42$ and $m = 13$, inflating the minimum
detectable response by a factor of $1.524$
(\texttt{results/inference/verification.json}, entry \texttt{E11}).

At the design of Section \ref{sec:sim} with $C = 5$ and $m = 13$,
\eqref{eq:mde} returns $\kappa^{\mathrm{Wald}}_{\mathrm{MDE}} = 0.363$ and
$\kappa^{\mathrm{dir}}_{\mathrm{MDE}} = 0.214$ for the linear wedge; solving the
same noncentrality equation at the nonlinear wedge that design actually
generates raises these to $0.450$ and $0.242$. The simulated magnitude at which
size-adjusted power first crosses $0.80$ in that design is $0.254$ for the
directed cluster sign-change test, $0.477$ for the directed label permutation
test and $0.566$ for the chi square Wald at its size-adjusted critical value.
Those crossings are interpolated on a coarse $\kappa$ grid with $600$
replications per cell and are accurate to roughly $0.05$. Sources:
\texttt{mde\_table.csv} for the closed forms, including $I_k$ and
$\lambda^*(0.05,0.20,13)$ above, and \texttt{perm\_mde\_vs\_sim.csv} for the
crossings.

The Wald statistic is not the right statistic against this alternative.

\begin{theorem}[optimal directed test]
\label{thm:umpi}
Consider the Gaussian limit experiment in which the category event studies
satisfy $b_c \sim N(\tau + \kappa d_c k, \Sigma)$ independently across $c$, with
$\tau \in \mathbb{R}^q$ a free nuisance and $\Sigma$ common across categories,
so that $\hat{\gamma}_P \sim N(\kappa k_P, V)$ exactly with $V =
\operatorname{Var}(\hat{\gamma}_P)$, and suppose $V$ and the response shape
$k_P$ are known. Among tests of $H_0: \kappa = 0$ against $H_1: \kappa > 0$ that
are invariant under the group of shifts $b_c \mapsto b_c + a$, $a \in
\mathbb{R}^q$, common to all categories, the test rejecting when $t = k_P'V^{-1}
\hat{\gamma}_P / \sqrt{k_P'V^{-1}k_P} > z_{1-\alpha}$ is uniformly most
powerful, with power $\Phi(\kappa \sqrt{k_P'V^{-1}k_P} - z_{1-\alpha})$.
\end{theorem}

Theorem \ref{thm:umpi} is stated for known $V$. In practice $V$ is unknown, and
the feasible directed statistic substitutes the CR1 estimate built from $G$
clusters; that is the statistic reported in Table \ref{tab:size}
and in the power table of the online Appendix. The substitution costs the normal calibration, and the feasible
statistic is the one that the label permutation of Theorem \ref{thm:exact} and
the cluster sign-change randomization of Theorem \ref{thm:signflip} calibrate
without a normal approximation.

Misspecifying the direction costs power and not level. Running the test at
$\tilde{k}$ when the truth is $\gamma_P = \kappa k_P$ makes $t =
\tilde{k}'V^{-1}\hat{\gamma}_P/\sqrt{\tilde{k}'V^{-1}\tilde{k}}$ normal with
unit variance and mean $\kappa \cos\vartheta \sqrt{k_P'V^{-1}k_P}$, where
$\vartheta$ is the angle between $\tilde{k}$ and $k_P$ in the inner product
$\langle a,b \rangle = a'V^{-1}b$. Its power is
$\Phi(\kappa\cos\vartheta\sqrt{k_P'V^{-1}k_P} - z_{1-\alpha})$, so the minimum
detectable response is the $\kappa^{\mathrm{dir}}_{\mathrm{MDE}}$ of
\eqref{eq:mde} divided by $\cos\vartheta$, and it is unbounded as $\tilde{k}$
approaches orthogonality to $k_P$ in that metric. The test is level $\alpha$ for
every choice of $\tilde{k}$.

The online Appendix works through which statistic the ratio describes, the
simulated consequences over $600$ replications per cell, the case of a
response restricted only to the nonnegative orthant, and how the ratio
depends on the calibration of each test.

\section{Simulations}
\label{sec:sim}

\paragraph{Design.} The simulated panel is a two-layer recording
model with count outcomes, described in full in the online Appendix and
summarized in the note to Table \ref{tab:size}. We vary $G \in
\{20,30,51,100\}$ with one unit per cluster, $C \in \{3,5,8\}$ categories,
$m = 13$ post-treatment restrictions and $T = 120$ periods, over $2{,}000$
replications per cell.

\begin{table}[!tbp]
\centering
\small
\begin{spacing}{1.0}
\caption{Null rejection rates, nominal level $0.05$}
\label{tab:size}
\begin{tabular}{llrrrr}
\toprule
 & & \multicolumn{4}{c}{Clusters $G$} \\
\cmidrule(lr){3-6}
Statistic and calibration & $C$ & 20 & 30 & 51 & 100 \\
\midrule
Wald, $\chi^2_m$                  & 3 & 0.920 & 0.693 & 0.393 & 0.172 \\
                                  & 5 & 0.905 & 0.686 & 0.389 & 0.187 \\
                                  & 8 & 0.917 & 0.690 & 0.356 & 0.172 \\
\addlinespace
$W/m$, $F(m,G-1)$                 & 3 & 0.838 & 0.586 & 0.311 & 0.143 \\
                                  & 5 & 0.821 & 0.558 & 0.309 & 0.153 \\
                                  & 8 & 0.834 & 0.571 & 0.283 & 0.138 \\
\addlinespace
Hotelling at $G$                  & 3 & 0.175 & 0.164 & 0.122 & 0.082 \\
                                  & 5 & 0.174 & 0.160 & 0.124 & 0.092 \\
                                  & 8 & 0.172 & 0.167 & 0.106 & 0.088 \\
\addlinespace
Hotelling at $G_{\mathrm{eff}}$   & 3 & 0.000 & 0.022 & 0.047 & 0.061 \\
                                  & 5 & 0.000 & 0.020 & 0.045 & 0.065 \\
                                  & 8 & 0.000 & 0.021 & 0.045 & 0.055 \\
\addlinespace
Label permutation, Wald           & 3 & 0.000 & 0.000 & 0.000 & 0.000 \\
                                  & 5 & 0.053 & 0.050 & 0.046 & 0.059 \\
                                  & 8 & 0.051 & 0.051 & 0.043 & 0.058 \\
\addlinespace
Label permutation, directed       & 3 & 0.000 & 0.000 & 0.000 & 0.000 \\
                                  & 5 & 0.052 & 0.053 & 0.048 & 0.052 \\
                                  & 8 & 0.046 & 0.050 & 0.043 & 0.051 \\
\addlinespace
Wild cluster bootstrap, Wald      & 3 & 0.053 & 0.055 & 0.053 & 0.061 \\
                                  & 5 & 0.064 & 0.062 & 0.051 & 0.067 \\
                                  & 8 & 0.061 & 0.061 & 0.051 & 0.056 \\
\addlinespace
Cluster sign change, Wald         & 3 & 0.027 & 0.036 & 0.046 & 0.058 \\
                                  & 5 & 0.022 & 0.041 & 0.050 & 0.066 \\
                                  & 8 & 0.033 & 0.041 & 0.048 & 0.053 \\
\addlinespace
Cluster sign change, directed     & 3 & 0.058 & 0.056 & 0.044 & 0.045 \\
                                  & 5 & 0.058 & 0.061 & 0.057 & 0.046 \\
                                  & 8 & 0.055 & 0.055 & 0.057 & 0.065 \\
\bottomrule
\end{tabular}

\smallskip
\begin{minipage}{0.95\textwidth}
\footnotesize Notes: rejection rates for the joint null $H_0: \gamma_P = 0$ of
\eqref{eq:null} in the pooled gradient regression \eqref{eq:pooled}, estimated
by OLS with the CR1 covariance clustered on $g$, at nominal level $0.05$.
Design: $m = 13$ post-treatment restrictions, $q = 18$ included horizons,
$T = 120$, adoption at $t = 48$, $72$ and $96$ by $20\%$ of units each with
$40\%$ never treated, one unit per cluster, exposure indices equally spaced on
$[0.1,1]$. The Hotelling rows refer $W$ to the $F$ law of \eqref{eq:Fexact}
evaluated at $G$ and at $G_{\mathrm{eff}}$, which is $0.727\,G$ at
$G = 20$, $30$ and $100$ and $0.714\,G$ at $G = 51$. Label permutation
randomizes the $C!$ assignments of the exposure index to the category data
arrays, enumerating all of them at $C = 3$ and $C = 5$ and using $998$ sampled
assignments together with the identity at $C = 8$. The wild cluster bootstrap
uses $499$ draws and the cluster sign-change test $999$, each including the
identity draw. The directed rows use the contrast $w = \mathbf{1}_m$.
$2{,}000$ replications per cell; the Monte Carlo standard error is $0.005$ at
a rate of $0.05$, $0.011$ at $0.39$ and $0.006$ at $0.92$. Sources:
\texttt{results/inference/size\_table.csv} and
\texttt{anatomy\_table.csv}.
\end{minipage}
\end{spacing}
\end{table}

\paragraph{Size.} Table \ref{tab:size} reports null rejection rates for
\eqref{eq:null} at nominal level $0.05$. The chi-square calibration rejects at
$0.920$, $0.693$, $0.393$ and $0.172$ at $G = 20$, $30$, $51$ and $100$ with
$C = 3$. Across $C$ at $G = 51$ the rates are $0.393$, $0.389$ and $0.356$ at
$C = 3$, $5$ and $8$, against $0.378$, $0.378$ and $0.378$ predicted by
Proposition \ref{prop:anatomy} evaluated at $G_{\mathrm{eff}}$, which moves in
$C$ only through $\omega_2$. The measured decline of $0.037$ from $C = 3$ to
$C = 8$ runs the other way. The two cells are independently seeded, so the
standard error of the difference is $0.015$ and the decline is $2.5$ standard
errors of it.
One candidate is the loading normalization, whose estimate improves with $C$,
and we have not isolated it.

The second block of rows refers $W/m$ to $F(m, G-1)$, the usual small-sample patch. Proposition \ref{prop:anatomy} puts the correct denominator degrees of
freedom at $G - m$, which is $38$ rather than $50$ at $G = 51$ and $m = 13$,
and the patch removes about a quarter of the excess over the nominal level:
$0.309$ against $0.389$ at $G = 51$, $C = 5$. The Hotelling calibration at $G$,
which is the one Proposition \ref{prop:anatomy} licenses under Assumption
\ref{as:I}, rejects between $0.082$ and $0.175$ over the twelve cells. That
calibration is model based rather than randomization based, and Assumption
\ref{as:I}(i) is what it rests on: with centered $\chi^2_1$ cluster
contributions, whose skewness is $2.83$, its size is $0.184$
(\texttt{verification.json}, entry \texttt{E4}). We report it as a
diagnostic and recommend the randomization calibrations of Sections
\ref{sec:perm} and \ref{sec:cluster} in its place.

The two-sided label permutation test is at nominal size for $C \in \{5,8\}$,
ranging from $0.043$ to $0.059$ over all eight such cells, and its directed
version from $0.043$ to $0.053$. At $C = 3$ both have rejection rate $0.000$,
which is Proposition \ref{prop:resolution}(iii) rather than a numerical
artifact. Three categories is not a rare design, since many recording systems
sort outcomes into a handful of coarse classes, and the resolution floor makes
label randomization unusable there at any conventional level however large the
panel. A researcher with $C = 3$ has three options: use the cluster level
procedures of Section \ref{sec:cluster}, which place no requirement on $C$ and
carry the sizes reported in the same table; split the outcome into finer
categories if the recording system records them separately, since the floor is
driven by $C!$ and not by sample size; or report the realized floor alongside
the $p$ value so that a nonrejection is not read as evidence. The third is a
statement about what the design can support. The wild cluster bootstrap ranges from $0.051$ to $0.067$. The
two-sided cluster sign-change test ranges from $0.022$ to $0.066$ and is
conservative at $G = 20$, where its three rates are $0.027$, $0.022$ and
$0.033$; the directed version rejects at $0.058$, $0.058$ and $0.055$ in those
same three cells.

\paragraph{The design rule.} Applying Corollary \ref{cor:asym} through
\eqref{eq:Fexact}, the largest number of restrictions holding chi-square size
below $0.075$ is $m = 2$ at $G_{\mathrm{eff}} = 36.4$, $m = 4$ at $72.7$, $m =
6$ at $100$ and $m = 8$ at $150$. Testing $m = 13$ restrictions at
$G_{\mathrm{eff}} = 36.4$ gives a predicted size of $0.379$. These are
evaluations of \eqref{eq:Fexact} at values of $m$ that no simulated cell runs;
every cell of every block in this section has $m = 13$, so the rule has no
simulated counterpart here.

\paragraph{Power.} Against a flat recording response of magnitude
$\kappa = 0.45$ at $G = 51$ and $C = 5$, the label permutation test with the two
sided Wald statistic attains size-adjusted power of $0.383$ while the cluster
sign-change test applied to the directed statistic attains $0.983$, which is the
ordering Corollary \ref{cor:mde} predicts. The online Appendix gives the power
table over the full grid of magnitudes, the contrast used, and the two cells at
$C = 5$ where the size-adjusted power of the two-sided permutation test is
nonmonotone in $\kappa$.

\paragraph{Further evidence.} The online Appendix reports seven further blocks of this design: the exact null law checked in an experiment built to satisfy Assumption \ref{as:I}, the accuracy of the substitution of $G_{\mathrm{eff}}$, the size distortion when category exchangeability fails and what studentizing does to it, the price label randomization charges relative to a normal benchmark, how the exposure spread, the panel length and the pre window move the minimum detectable response, the behavior of the calibrations in four harder designs with clustered, unequal, serially correlated and heavy tailed errors, and a benchmark against the bias-reduced variance of \citet{bell2002bias}. None of them changes a conclusion drawn here.


\section{What the design has to satisfy}
\label{sec:designcalc}

A falsification test is worth running only against responses it could detect,
and whether it can is a property of the design that is computable before the
outcome is examined. Write $\kappa_{\mathrm{MDE}}$ for the response the design
detects with probability $0.80$ at level $\alpha$, which by Corollary
\ref{cor:mde} depends on the number of restrictions, the path information
$I_k$, the exposure spread $\Sigma_\delta^2$ and the residual variance.

What $\kappa_{\mathrm{MDE}}$ has to be compared with needs care, because the
test does not see the whole of a recording response. Let $m_c$ be the log
response the recording change produces in category $c$. The alternative the
gradient test is built against has cross-category shape proportional to
$\delta_c$, so the part of $m$ that the test can detect is its projection on the
exposure index,
\begin{equation}
\kappa^{\mathrm{inst}} = \frac{\sum_c \delta_c m_c}{\sum_c \delta_c^2} ,
\label{eq:kinst}
\end{equation}
and the design is adequate when $\kappa_{\mathrm{MDE}} \le
\kappa^{\mathrm{inst}}$. Both inputs are external to the outcome panel, $d_c$
from the recording institution and $m_c$ from whatever is known about the
recording change, so \eqref{eq:kinst} is computable in advance, which is the
only reason the comparison is worth making. A quantity read off the
post-treatment estimates, a per category breakdown value for instance, is not a
candidate for $\kappa^{\mathrm{inst}}$: it would make the screen a function of
the outcome it screens.

Five features of the specification move the ratio.

The statistic is the first. By Corollary \ref{cor:mde} the directed statistic
has the smaller minimum detectable response whenever the response path is
signed, which the recording mechanism makes it. The pre window is the second,
and Corollary \ref{cor:blind} is what makes narrowing it free: the pre-period carries no information about a recording response, so horizons spent there buy
nothing this test can use while competing with the post block for
identification. Dropping the crime pre window from six horizons to one raises
$I_k$ from $40.5$ to $175.5$ and cuts $\kappa_{\mathrm{MDE}}$ by the square root
of that. The number of restrictions is the third, through the critical value.
The fourth is the exposure spread, and the cheapest way to widen it is a
category the recording system has no discretion over, which anchors the low end
of $d$ at zero.

The fifth is alignment, and it is the one that can defeat the other four. By
\eqref{eq:kinst} the numerator is an inner product, so a design can have ample
power against a response of the average size the institution implies and almost
none against the response it actually implies, if the index orders the
categories on a different principle from the one the recording system acts on.
The cosine between $\delta$ and $m$ is the whole of it, and it costs one inner
product to check. Section \ref{sec:app} reports a design that fails here while
passing on every other lever, which is the case a researcher most needs to be
warned about, because widening the spread and narrowing the pre window both
look like progress and neither helps.

\section{Application}
\label{sec:app}

We apply the procedures to two state panels assembled for the study of
treatment-induced recording change. The identification claim in each case is
that under measurement invariance the post-treatment event-study coefficients do
not vary with the exposure index, so \eqref{eq:null} holds.

The first four levers of Section \ref{sec:designcalc} move the crime panel a
long way. Homicide, which the panel already carries and which the specification
as originally run holds out as a $d = 0$ series, more than doubles
$\Sigma_\delta^2$ from $0.349$ to $0.769$ when it enters the gradient; dropping
the pre window from six horizons to one raises $I_k$ from $40.5$ to $175.5$; and
the directed statistic takes the rest. Over the forty eight specifications we
enumerate the minimum detectable response falls to $0.052$ at its best, with
$C = 7$, $m = 9$ and $G_{\mathrm{eff}} = 42.9$.

What it has to beat is far smaller. The FBI's comparison of the two reporting
systems over $6{,}299$ agencies puts the log change in recorded counts at under
$0.001$ for rape, $0.006$ for robbery and for aggravated assault, $0.010$ for
burglary, $0.026$ for larceny and $0.027$ for motor vehicle theft, with murder
unchanged, averaging $0.011$ across the seven index categories (Congressional
Research Service report R46668, which also reports the earlier Bureau of Justice
Statistics comparison over $1{,}131$ agencies). Projected on the exposure index
through \eqref{eq:kinst} the first gives $\kappa^{\mathrm{inst}} = 0.0044$ and
the second $0.0017$. No specification of the crime panel is adequate, the best
ratio being $11.8$.

The shortfall is misalignment rather than power, and that distinction is the
useful part. The responses themselves are of the order the design detects; it is
their projection on $\delta$ that is not, the correlation between the two being
$0.14$ under the first comparison and $0.03$ under the second. The exposure index
grades categories by the share of victimizations not reported to police, and the
move to incident-based reporting does not act on victim reporting at all. It
drops the hierarchy rule and makes co-occurring offenses countable, which is a
change in police counting, and the two orderings come apart and in places
reverse: motor vehicle theft has the second lowest $d_c$ and nearly the largest
response, because vehicle theft often accompanies another index offense, while
rape has nearly the highest $d_c$ and no response worth measuring, because it
rarely accompanies a more serious one. Widening the spread and narrowing the pre
window both looked like progress here and neither helped.

The same calculation screens the health panels out, the best ratios being $12.5$
in the annual sub design and $14.8$ in the odd-year sub design against $33.9$
and $53.2$ as run. The exposure index there already places three conditions at
$d = 0$, so widening the spread has little left to buy, and the shortfall is
again an order of magnitude.

We therefore report all three designs as illustrations of the calculus screening
a design out, and none of them as evidence about incident-based reporting or
about Medicaid expansion. Across every specification of both panels the
randomization calibrations reject at about their nominal rate, ten of the two
hundred and forty three crime $p$ values and nineteen of the three hundred and
seventy eight health ones falling at or below $0.05$, which is what the null
predicts and what a design this far from adequate can be expected to produce
either way. Those nonrejections carry no information about measurement
invariance. The point of computing $\kappa_{\mathrm{MDE}}$ and
$\kappa^{\mathrm{inst}}$ before the outcome is that all of this was knowable in
advance. Sources: \texttt{design\_power.csv} for the screen and
\texttt{theory/institutional\_wedge.md} for the derivation of
$\kappa^{\mathrm{inst}}$.

The online Appendix describes the two panels, the loading normalization
applied to each category array, and what the chi-square calibration returns
on them in contrast to the randomization calibrations.

\paragraph{Data availability.} The recorded crime series are the FBI's
summary based estimated state totals; the NIBRS adoption dates come from the FBI
Crime Data Explorer agency API; the reporting rates defining $d_c$ come from
Bureau of Justice Statistics, \emph{Criminal Victimization 2019}, NCJ 255113,
Table 6. The health series are BRFSS crude prevalence from the CDC Socrata
dataset \texttt{dttw-5yxu}, and the expansion effective dates come from the
Kaiser Family Foundation table of state Medicaid expansion decisions. All are
public. The two constructed panels and the code that builds them are in the
replication archive, whose location is withheld for anonymous review, under
\texttt{data/processed/} and \texttt{scripts/}. The per category changes in
recorded counts underlying $\kappa^{\mathrm{inst}}$ are from Congressional
Research Service report R46668, which summarizes the FBI and Bureau of Justice
Statistics comparisons of the two reporting systems; the projection is derived
in \texttt{theory/institutional\_wedge.md}. Every $p$ value,
$G_{\mathrm{eff}}$ and predicted size is read from
\texttt{results/empirics/inference\_gradient.csv} and every realized resolution
floor from \texttt{results/empirics/app\_floors\_injection.csv}. The forty
eight screened specifications of each panel, with their minimum detectable
responses and adequacy ratios, are in
\texttt{results/inference/design\_power.csv}.

\section{Conclusion}

We have argued that a treatment-induced change in recording cannot be separated
from a treatment effect by anything computed from the outcome panel. Under
Assumption \ref{as:rec} the deviation that reproduces such a change is exactly
zero before treatment, so no diagnostic of that kind rejects more often than its
nominal level, and the relative magnitudes class scales a pre-period quantity
that is identically zero and therefore admits no deviation at any value of
$\bar M$. Its coverage of the latent effect falls from $0.955$ to $0.000$ as the panel
grows, while a bound set from the recording institution instead of from the pre-period holds at $1.000$ throughout.

Information about such a change must come from the cross section of outcome
categories, and the test that uses it needs a reference distribution of its own.
We have shown that the chi-square distribution is badly oversized at the cluster
counts these tests are run on, since at $G = 51$ and $m = 13$ its size is
$0.389$. We therefore give three alternatives and characterize each. The label
permutation test of Definition \ref{def:perm} is exact in finite samples for
\eqref{eq:null} under Assumption \ref{as:CE} and requires no condition on $G$,
the cluster sign-change test is exact under sign symmetry of the cluster
contributions, and the bootstrap carries no exactness claim. Their simulated
sizes in Table \ref{tab:size} are $0.043$--$0.059$, $0.022$--$0.066$ and
$0.051$--$0.067$ at $C \in \{5,8\}$.

Three limitations bound what we claim. Proposition \ref{prop:rm} is proved only
for a single adoption cohort, where the online Appendix shows by simulation that
the conclusion survives staggering. Proposition \ref{prop:anatomy} requires equal
design weights, which no design we study satisfies, so we use it to explain the
distortion rather than to replace it. Assumption \ref{as:CE} fails when categories
have different error variances: in a Gaussian limit experiment that failure is
liberal, at size $0.102$ for a variance ratio of $25$ at $C = 5$, while in the
panel, where the CR1 covariance is recomputed for every permuted label vector, it
is conservative instead and costs power. The cluster level procedures do not rely
on the assumption at all.

Lemma \ref{lem:sep} requires a category balanced design, a restriction the
health panel shows to be binding, and what blocks Theorem \ref{thm:exact} there
is Assumption \ref{as:CE}, because the pattern of observed cells depends on the
category label. Extending the reduction to unbalanced panels would recover the
procedure in settings where categories are surveyed on different schedules, and
we leave that extension to future work.

\section*{Statements and Declarations}

\paragraph{Funding.} No funding was received for conducting this study.

\paragraph{Competing interests.} The author has no competing interests to
declare that are relevant to the content of this article.

\paragraph{Ethics approval.} Not applicable. The study uses only aggregated
public use statistics and involves no human participants, their data or
biological material.

\paragraph{Data, material and code availability.} Every series used here is a
public use file. The sources, the two constructed panels and the code that
builds them are listed under data availability in Section \ref{sec:app},
together with the replication archive that holds them and the output file
behind every number reported in this paper.

\paragraph{Use of generative AI.} The author used a generative AI assistant
during preparation of this manuscript, for drafting and editing prose and
formatting the reference list, and for some of the simulation code. This goes
beyond AI assisted copy editing and is declared here for that reason. All
theoretical results and their proofs, the design of the empirical analysis, and
the interpretation of all findings are the author's own. The author has verified
every reported number against the output files named in the notes to the tables
and takes full responsibility for the content.

\paragraph{Online Resource.} The online Appendix referred to throughout
contains every proof and the supporting simulation and application evidence. It
restates nothing and follows the numbering used here, and is submitted as
Online Resource 1.

\begin{thebibliography}{99}

\bibitem[Armstrong and Koles\'ar(2020)]{armstrong2020simple} Armstrong, T.~B., and Koles\'ar, M. (2020), ``Simple and Honest Confidence Intervals in Nonparametric Regression,'' \emph{Quantitative Economics}, 11, 1--39, DOI: 10.3982/qe1199.

\bibitem[Arora and Chand(2026)]{arora2026trading} Arora, P., and Chand, A. (2026), ``Trading Scope for Credibility in Difference-in-differences,'', arXiv:2608.16867.

\bibitem[Baicker et al.(2013)]{baicker2013oregon} Baicker, K., Taubman, S.~L., Allen, H.~L., Bernstein, M., Gruber, J.~H., Newhouse, J.~P., Schneider, E.~C., Wright, B.~J., Zaslavsky, A.~M., and Finkelstein, A.~N. (2013), ``The Oregon Experiment: Effects of Medicaid on Clinical Outcomes,'' \emph{New England Journal of Medicine}, 368, 1713--1722, DOI: 10.1056/nejmsa1212321.

\bibitem[Bell and McCaffrey(2002)]{bell2002bias} Bell, R.~M., and McCaffrey, D.~F. (2002), ``Bias Reduction in Standard Errors for Linear Regression With Multi-stage Samples,'' \emph{Survey Methodology}, 28, 169--181.

\bibitem[Bester et al.(2011)]{bester2011cluster} Bester, C.~A., Conley, T.~G., and Hansen, C.~B. (2011), ``Inference With Dependent Data Using Cluster Covariance Estimators,'' \emph{Journal of Econometrics}, 165, 137--151, DOI: 10.1016/j.jeconom.2011.01.007.

\bibitem[Borusyak et al.(2024)]{borusyak2024revisiting} Borusyak, K., Jaravel, X., and Spiess, J. (2024), ``Revisiting Event-study Designs: Robust and Efficient Estimation,'' \emph{Review of Economic Studies}, 91, 3253--3285, DOI: 10.1093/restud/rdae007.

\bibitem[Callaway and Sant'Anna(2021)]{callaway2021did} Callaway, B., and Sant'Anna, P.~H.~C. (2021), ``Difference-in-differences With Multiple Time Periods,'' \emph{Journal of Econometrics}, 225, 200--230, DOI: 10.1016/j.jeconom.2020.12.001.

\bibitem[Cai and Low(2004)]{cailow2004adaptation} Cai, T.~T., and Low, M.~G. (2004), ``An Adaptation Theory for Nonparametric Confidence Intervals,'' \emph{The Annals of Statistics}, 32, 1805--1840, DOI: 10.1214/009053604000000049.

\bibitem[Cameron et al.(2008)]{cameron2008bootstrap} Cameron, A.~C., Gelbach, J.~B., and Miller, D.~L. (2008), ``Bootstrap-based Improvements for Inference With Clustered Errors,'' \emph{Review of Economics and Statistics}, 90, 414--427, DOI: 10.1162/rest.90.3.414.

\bibitem[Campbell(1969)]{campbell1969reforms} Campbell, D.~T. (1969), ``Reforms as Experiments,'' \emph{American Psychologist}, 24, 409--429, DOI: 10.1037/h0027982.

\bibitem[Campbell and Stanley(1963)]{campbell1963experimental} Campbell, D.~T., and Stanley, J.~C. (1963), \emph{Experimental and Quasi-Experimental Designs for Research}, Chicago: Rand McNally.

\bibitem[Carter et al.(2017)]{carter2017cluster} Carter, A.~V., Schnepel, K.~T., and Steigerwald, D.~G. (2017), ``Asymptotic Behavior of a $t$-test Robust to Cluster Heterogeneity,'' \emph{Review of Economics and Statistics}, 99, 698--709, DOI: 10.1162/rest\_a\_00639.

\bibitem[de Chaisemartin and D'Haultf\oe uille(2020)]{dechaisemartin2020twoway} de Chaisemartin, C., and D'Haultf\oe uille, X. (2020), ``Two-way Fixed Effects Estimators With Heterogeneous Treatment Effects,'' \emph{American Economic Review}, 110, 2964--2996, DOI: 10.1257/aer.20181169.

\bibitem[Djogbenou et al.(2019)]{djogbenou2019asymptotic} Djogbenou, A.~A., MacKinnon, J.~G., and Nielsen, M.~{\O}. (2019), ``Asymptotic Theory and Wild Bootstrap Inference With Clustered Errors,'' \emph{Journal of Econometrics}, 212, 393--412, DOI: 10.1016/j.jeconom.2019.04.035.

\bibitem[Freeman(2026)]{freeman2026local} Freeman, H. (2026), ``Local Asymptotic Power of Honest Confidence Intervals,'', arXiv:2607.10558.

\bibitem[Geruso and Layton(2020)]{geruso2020upcoding} Geruso, M., and Layton, T. (2020), ``Upcoding: Evidence From Medicare on Squishy Risk Adjustment,'' \emph{Journal of Political Economy}, 128, 984--1026, DOI: 10.1086/704756.

\bibitem[Harasztosi and Lindner(2019)]{harasztosi2019minwage} Harasztosi, P., and Lindner, A. (2019), ``Who Pays for the Minimum Wage?'' \emph{American Economic Review}, 109, 2693--2727, DOI: 10.1257/aer.20171445.

\bibitem[Hotelling(1931)]{hotelling1931} Hotelling, H. (1931), ``The Generalization of Student's Ratio,'' \emph{Annals of Mathematical Statistics}, 2, 360--378.

\bibitem[Iyer et al.(2012)]{iyer2012political} Iyer, L., Mani, A., Mishra, P., and Topalova, P. (2012), ``The Power of Political Voice: Women's Political Representation and Crime in India,'' \emph{American Economic Journal: Applied Economics}, 4, 165--193, DOI: 10.1257/app.4.4.165.

\bibitem[Lehmann and Romano(2005)]{lehmann2005testing} Lehmann, E.~L., and Romano, J.~P. (2005), \emph{Testing Statistical Hypotheses}, (3rd ed.). New York: Springer.

\bibitem[Lipsitch et al.(2010)]{lipsitch2010negative} Lipsitch, M., Tchetgen Tchetgen, E., and Cohen, T. (2010), ``Negative Controls: A Tool for Detecting Confounding and Bias in Observational Studies,'' \emph{Epidemiology}, 21, 383--388, DOI: 10.1097/ede.0b013e3181e4bfd7.

\bibitem[Low(1997)]{low1997nonparametric} Low, M.~G. (1997), ``On Nonparametric Confidence Intervals,'' \emph{The Annals of Statistics}, 25, 2547--2554.

\bibitem[MacKinnon and Webb(2017)]{mackinnon2017wild} MacKinnon, J.~G., and Webb, M.~D. (2017), ``Wild Bootstrap Inference for Wildly Different Cluster Sizes,'' \emph{Journal of Applied Econometrics}, 32, 233--254, DOI: 10.1002/jae.2508.

\bibitem[Miller and Segal(2019)]{miller2019female} Miller, A.~R., and Segal, C. (2019), ``Do Female Officers Improve Law Enforcement Quality? Effects on Crime Reporting and Domestic Violence,'' \emph{Review of Economic Studies}, 86, 2220--2247, DOI: 10.1093/restud/rdy051.

\bibitem[N\'u\~nez and Matthews(2025)]{nunez2025detection} N\'u\~nez, I., and Matthews, A.~A. (2025), ``Detection Bias and the Role of Negative Control Outcomes,'' \emph{BMJ Medicine}, 4, e001336, DOI: 10.1136/bmjmed-2025-001336.

\bibitem[Ody et al.(2019)]{ody2019readmissions} Ody, C., Msall, L., Dafny, L.~S., Grabowski, D.~C., and Cutler, D.~M. (2019), ``Have Hospital Readmissions Declined in the Wake of the Hospital Readmissions Reduction Program?'' \emph{Health Affairs}, 38, 36--43.

\bibitem[Pustejovsky and Tipton(2018)]{pustejovsky2018small} Pustejovsky, J.~E., and Tipton, E. (2018), ``Small-sample Methods for Cluster-robust Variance Estimation and Hypothesis Testing in Fixed Effects Models,'' \emph{Journal of Business \& Economic Statistics}, 36, 672--683, DOI: 10.1080/07350015.2016.1247004.

\bibitem[Rambachan and Roth(2023)]{rambachan2023credible} Rambachan, A., and Roth, J. (2023), ``A More Credible Approach to Parallel Trends,'' \emph{Review of Economic Studies}, 90, 2555--2591, DOI: 10.1093/restud/rdad018.

\bibitem[Roth(2022)]{roth2022pretest} Roth, J. (2022), ``Pretest With Caution: Event-study Estimates After Testing for Parallel Trends,'' \emph{American Economic Review: Insights}, 4, 305--322, DOI: 10.1257/aeri.20210236.

\bibitem[Roth et al.(2023)]{roth2023trending} Roth, J., Sant'Anna, P.~H.~C., Bilinski, A., and Poe, J. (2023), ``What's Trending in Difference-in-differences? a Synthesis of the Recent Econometrics Literature,'' \emph{Journal of Econometrics}, 235, 2218--2244, DOI: 10.1016/j.jeconom.2023.03.008.

\bibitem[Singh(2020)]{singh2020myths} Singh, A. (2020), ``Myths of Official Measurement: Auditing and Improving Administrative Data in Developing Countries,'' RISE Working Paper 20/042.

\bibitem[Singh(2024)]{singh2024administrative} Singh, A. (2024), ``Improving Administrative Data at Scale: Experimental Evidence on Digital Testing in Indian Schools,'' \emph{The Economic Journal}, 134, 2207--2223, DOI: 10.1093/ej/ueae017.

\bibitem[Sun and Abraham(2021)]{sunabraham2021} Sun, L., and Abraham, S. (2021), ``Estimating Dynamic Treatment Effects in Event Studies With Heterogeneous Treatment Effects,'' \emph{Journal of Econometrics}, 225, 175--199, DOI: 10.1016/j.jeconom.2020.09.006.

\bibitem[Tipton and Pustejovsky(2015)]{tipton2015small} Tipton, E., and Pustejovsky, J.~E. (2015), ``Small-sample Adjustments for Tests of Moderators and Model Fit Using Robust Variance Estimation in Meta-regression,'' \emph{Journal of Educational and Behavioral Statistics}, 40, 604--634, DOI: 10.3102/1076998615606099.

\bibitem[Weisburd et al.(2024)]{weisburd2024crs} Weisburd, D., Wilson, D.~B., Gill, C., Kuen, K., and Zastrow, T. (2024), ``Are We Underestimating the Crime Prevention Outcomes of Community Policing? the Importance of Crime Reporting Sensitivity Bias,'' \emph{Journal of Law and Empirical Analysis}, 1, 1--17, DOI: 10.1177/2755323x241233469.

\bibitem[Zhang et al.(2025)]{zhang2025nco} Zhang, D., Zhang, B., Wang, H., Lu, Y., Wolock, C.~J., Hu, W., Wang, L., Hripcsak, G., and Chen, Y. (2025), ``Negative Control-calibrated Difference-in-difference Analyses: Addressing Unmeasured Confounding in RWD With Application to Racial/ethnic Differences,'' \emph{npj Digital Medicine}, 8, 452, DOI: 10.1038/s41746-025-01821-w.

\end{thebibliography}

\end{document}
